% !Mode:: "TeX:UTF-8"
% !TEX root = ../../main.tex
% ============================================================
% TOPIC 19
% ============================================================
\mashgulot{19}{Eigenvalues and eigenvectors of a linear operator}

\begin{tarif}
A number $\lambda$ is called an \textbf{eigenvalue} of the operator
$\mathcal A$ (or of its matrix $A$) if there is a nonzero vector $x$ such
that $\mathcal Ax=\lambda x$; every such vector $x$ is called an
\textbf{eigenvector} corresponding to $\lambda$.
\end{tarif}

The equation $Ax=\lambda x$ is equivalent to the homogeneous system
$(A-\lambda E)x=0$. This system has a nonzero solution if and only if its
matrix is singular, that is,
\begin{equation}
\det(A-\lambda E)=0 .
\end{equation}
Equation (\theequation) is called the \textbf{characteristic equation}
of $A$, and its left-hand side, a polynomial of degree $n$ in $\lambda$,
is the \textbf{characteristic polynomial}.

The computation therefore splits into two stages. We first solve the
characteristic equation and obtain the eigenvalues
$\lambda_1,\dots,\lambda_n$. Then, for each eigenvalue $\lambda_k$, we
solve the homogeneous system $(A-\lambda_kE)x=0$; a fundamental system of
its solutions is a maximal set of linearly independent eigenvectors
belonging to $\lambda_k$.

\textbf{Properties.} The eigenvalues, counted with multiplicity, satisfy
$\lambda_1+\dots+\lambda_n=\operatorname{tr}A$ and
$\lambda_1\cdots\lambda_n=\det A$. The eigenvalues of a triangular matrix
are its diagonal entries. If $\lambda$ is an eigenvalue of $A$, then
$\lambda^k$ is an eigenvalue of $A^k$ and, for nonsingular $A$,
$1/\lambda$ is an eigenvalue of $A^{-1}$, with the same eigenvector.
Eigenvectors belonging to distinct eigenvalues are linearly independent.
Finally, a real symmetric matrix has only real eigenvalues, and its
eigenvectors belonging to distinct eigenvalues are orthogonal.

\begin{misol}
Find the eigenvalues and eigenvectors of
$A=\begin{pmatrix}4&1\\2&3\end{pmatrix}$.

\textbf{Solution.} The characteristic equation is
\[\kf{\det(A-\lambda E)=\begin{vmatrix}4-\lambda&1\\2&3-\lambda\end{vmatrix}
=\lambda^2-7\lambda+10=0,}\]
so $\lambda_1=2$ and $\lambda_2=5$. For $\lambda_1=2$ the system
$(A-2E)x=0$ reduces to the single equation $2x_1+x_2=0$, whence
$x^{(1)}=c(1;-2)$. For $\lambda_2=5$ we obtain $-x_1+x_2=0$ (the second
equation $2x_1-2x_2=0$ is proportional to it), so $x^{(2)}=c(1;1)$. In
both cases $c\neq0$.

As a check, $A(1;1)^T=(5;5)^T=5(1;1)^T$; moreover,
$\lambda_1+\lambda_2=7=\operatorname{tr}A$ and
$\lambda_1\lambda_2=10=\det A$.
\end{misol}

\begin{misol}
Find the eigenvalues and eigenvectors of
$A=\begin{pmatrix}2&0&0\\0&3&4\\0&4&-3\end{pmatrix}$.

\textbf{Solution.} Expanding along the first row, we get
\[\kf{\det(A-\lambda E)=(2-\lambda)\big[(3-\lambda)(-3-\lambda)-16\big]
=(2-\lambda)(\lambda^2-25)=0,}\]
so $\lambda_1=2$, $\lambda_2=5$, $\lambda_3=-5$. For $\lambda=2$ the
system gives $x_2=x_3=0$ with $x_1$ arbitrary, hence $x^{(1)}=(1;0;0)$.
For $\lambda=5$ we have $x_1=0$ and $-2x_2+4x_3=0$, hence
$x^{(2)}=(0;2;1)$. For $\lambda=-5$ we have $x_1=0$ and $8x_2+4x_3=0$,
hence $x^{(3)}=(0;1;-2)$.

Since $A$ is symmetric, the three eigenvectors must be pairwise
orthogonal; indeed, $(0;2;1)\cdot(0;1;-2)=0$, and both vectors are
clearly orthogonal to $(1;0;0)$.
\end{misol}

\begin{misol}
For the matrix $A=\begin{pmatrix}0&-1\\1&0\end{pmatrix}$ of the rotation
of the plane through $90^\circ$, the characteristic equation is
$\lambda^2+1=0$, which has no real roots. This agrees with geometry: a
rotation through $90^\circ$ leaves no nonzero vector on its own line.
\end{misol}

\subsection{Applications of eigenvalues and eigenvectors in information security}

\textbf{PageRank and reputation systems.} The importance of the nodes of a
network (web sites, users) is computed as the eigenvector of the
transition matrix corresponding to the eigenvalue $\lambda=1$. Trust
scores used to detect phishing sites and botnets are obtained in exactly
this way.

\begin{misol}
For the transition matrix
$P=\begin{pmatrix}0.9&0.2\\0.1&0.8\end{pmatrix}$ the system $(P-E)x=0$
reduces to $-0.1x_1+0.2x_2=0$, that is, $x_1=2x_2$. Normalizing so that
the components add up to one, we obtain the stationary distribution
$x=\left(\frac23;\frac13\right)$: in the long run the system spends
two thirds of the time in state~1.
\end{misol}

\textbf{Spectral analysis and anomaly detection.} The eigenvectors of the
covariance matrix of network traffic that belong to its largest
eigenvalues (the principal components of PCA) span the subspace of
“normal” traffic; observations lying far from this subspace are flagged
as possible attacks. The eigenvalues of the adjacency matrix of a graph,
in turn, describe the connectivity of a network and reveal its weak
points.

\subsection{Computing with Python, Maple and Mathcad}

All three systems find eigenvalues and eigenvectors directly. We take the
matrix $A=\begin{pmatrix}4&1\\2&3\end{pmatrix}$ of the first example. In
Python, SymPy factors the characteristic polynomial and returns, for each
eigenvalue, its multiplicity and a basis of eigenvectors.
\begin{Verbatim}[fontsize=\footnotesize,frame=single,label={\small Python}]
import sympy as sp
A = sp.Matrix([[4, 1], [2, 3]])
lam = sp.symbols('lambda')
print(sp.factor((A - lam*sp.eye(2)).det()))
for val, mult, vecs in A.eigenvects():
    print(val, mult, list(vecs[0]))
\end{Verbatim}
\begin{Verbatim}[fontsize=\footnotesize,frame=single,label={\small Output}]
(lambda - 5)*(lambda - 2)
2 1 [-1/2, 1]
5 1 [1, 1]
\end{Verbatim}
The eigenvector $(-\frac12;1)$ is a multiple of $(1;-2)$, as it should be.

In Maple the same computation uses the \texttt{LinearAlgebra} package;
\texttt{Eigenvectors} returns the column of eigenvalues and a matrix
whose columns are the corresponding eigenvectors.
\begin{Verbatim}[fontsize=\footnotesize,frame=single,label={\small Maple}]
with(LinearAlgebra):
A := Matrix([[4, 1], [2, 3]]):
CharacteristicPolynomial(A, lambda);
                 lambda^2 - 7 lambda + 10
Eigenvectors(A);
                 [2]  [-1/2  1]
                 [5], [ 1    1]
\end{Verbatim}

Mathcad Prime has the built-in functions \texttt{eigenvals(A)},
\texttt{eigenvecs(A)} and \texttt{eigenvec(A, z)}; the eigenvectors are
returned normalized to unit length.
\begin{Verbatim}[fontsize=\footnotesize,frame=single,label={\small Mathcad}]
A := | 4  1 |
     | 2  3 |
eigenvals(A) = | 2 |
               | 5 |
eigenvec(A, 2) = | -0.447 |     eigenvec(A, 5) = | 0.707 |
                 |  0.894 |                      | 0.707 |
\end{Verbatim}

\subsection{Exercises on eigenvalues and eigenvectors}

\emph{1. Find the eigenvalues and eigenvectors of the following
$2\times2$ matrices.}
\bigskip
\begin{multicols}{2}\raggedright
\begin{enumerate}[label=\arabic*., leftmargin=*]
\item $\begin{pmatrix}4 & 1\\ 0 & 5\end{pmatrix}$
\item $\begin{pmatrix}0 & -2\\ 0 & 2\end{pmatrix}$
\item $\begin{pmatrix}-4 & -4\\ 0 & 0\end{pmatrix}$
\item $\begin{pmatrix}-9 & -6\\ 12 & 9\end{pmatrix}$
\item $\begin{pmatrix}10 & 5\\ -10 & -5\end{pmatrix}$
\item $\begin{pmatrix}0 & 8\\ 0 & 4\end{pmatrix}$
\item $\begin{pmatrix}5 & 4\\ -2 & -1\end{pmatrix}$
\item $\begin{pmatrix}4 & 0\\ 6 & -2\end{pmatrix}$
\item $\begin{pmatrix}2 & 0\\ 4 & 4\end{pmatrix}$
\item $\begin{pmatrix}-2 & 2\\ 0 & -3\end{pmatrix}$
\item $\begin{pmatrix}1 & 0\\ 2 & 0\end{pmatrix}$
\item $\begin{pmatrix}5 & 2\\ -1 & 2\end{pmatrix}$
\item $\begin{pmatrix}1 & 0\\ -3 & -2\end{pmatrix}$
\item $\begin{pmatrix}0 & -4\\ 0 & -2\end{pmatrix}$
\item $\begin{pmatrix}5 & 0\\ 3 & 2\end{pmatrix}$
\item $\begin{pmatrix}-4 & -4\\ 0 & -2\end{pmatrix}$
\item $\begin{pmatrix}-1 & -6\\ 0 & 5\end{pmatrix}$
\item $\begin{pmatrix}3 & 6\\ 0 & 0\end{pmatrix}$
\item $\begin{pmatrix}5 & 0\\ -7 & -2\end{pmatrix}$
\item $\begin{pmatrix}-1 & 0\\ 12 & 5\end{pmatrix}$
\item $\begin{pmatrix}1 & 0\\ -4 & 3\end{pmatrix}$
\item $\begin{pmatrix}2 & -2\\ 0 & 1\end{pmatrix}$
\item $\begin{pmatrix}3 & 3\\ -6 & -6\end{pmatrix}$
\item $\begin{pmatrix}5 & 8\\ 0 & 1\end{pmatrix}$
\item $\begin{pmatrix}-3 & 8\\ 0 & 1\end{pmatrix}$
\item $\begin{pmatrix}-2 & 1\\ 0 & -1\end{pmatrix}$
\item $\begin{pmatrix}3 & -4\\ 2 & -3\end{pmatrix}$
\item $\begin{pmatrix}-1 & 2\\ 0 & 1\end{pmatrix}$
\item $\begin{pmatrix}5 & 0\\ 10 & 0\end{pmatrix}$
\item $\begin{pmatrix}-4 & -4\\ 0 & 0\end{pmatrix}$
\end{enumerate}
\end{multicols}

\emph{2. Find the eigenvalues and eigenvectors of the following
$3\times3$ matrices.}
\bigskip
\begin{enumerate}[label=\arabic*., leftmargin=*]
\item $\begin{pmatrix}-8 & 0 & -6\\ -8 & -1 & -9\\ 12 & 0 & 10\end{pmatrix}$
\item $\begin{pmatrix}-2 & -5 & -10\\ 0 & 1 & -2\\ 0 & 1 & 4\end{pmatrix}$
\item $\begin{pmatrix}3 & 0 & 0\\ 0 & -2 & 0\\ 1 & 6 & 4\end{pmatrix}$
\item $\begin{pmatrix}0 & 0 & 0\\ 2 & -2 & 6\\ 4 & 0 & 4\end{pmatrix}$
\item $\begin{pmatrix}1 & 0 & 0\\ 6 & -2 & 0\\ -6 & 2 & -1\end{pmatrix}$
\item $\begin{pmatrix}-3 & -2 & 2\\ 0 & -1 & -2\\ -4 & -4 & 1\end{pmatrix}$
\item $\begin{pmatrix}-1 & -2 & -2\\ 2 & -2 & -2\\ -2 & 0 & 0\end{pmatrix}$
\item $\begin{pmatrix}8 & 0 & -5\\ 0 & -1 & 0\\ 10 & 0 & -7\end{pmatrix}$
\item $\begin{pmatrix}-2 & -12 & -12\\ 0 & 4 & 3\\ 0 & 0 & 1\end{pmatrix}$
\item $\begin{pmatrix}-2 & 10 & -12\\ 0 & 3 & -6\\ 0 & 0 & -3\end{pmatrix}$
\item $\begin{pmatrix}0 & -9 & 1\\ 0 & 4 & 0\\ 0 & 3 & 1\end{pmatrix}$
\item $\begin{pmatrix}5 & 2 & 2\\ -6 & -2 & -4\\ 2 & 1 & 3\end{pmatrix}$
\item $\begin{pmatrix}-3 & 5 & -8\\ 0 & 2 & -8\\ 0 & 0 & -2\end{pmatrix}$
\item $\begin{pmatrix}-9 & 12 & 0\\ -6 & 9 & 0\\ 8 & -6 & 2\end{pmatrix}$
\item $\begin{pmatrix}2 & 1 & 6\\ 4 & -1 & 6\\ 1 & -1 & -3\end{pmatrix}$
\item $\begin{pmatrix}8 & 0 & 4\\ 5 & 1 & 2\\ -8 & 0 & -4\end{pmatrix}$
\item $\begin{pmatrix}-4 & 8 & -6\\ -2 & 4 & -4\\ -1 & 2 & -3\end{pmatrix}$
\item $\begin{pmatrix}-2 & 0 & 0\\ 4 & 2 & 0\\ -2 & 4 & 4\end{pmatrix}$
\item $\begin{pmatrix}-3 & 5 & 5\\ 2 & 0 & 2\\ -2 & 2 & 0\end{pmatrix}$
\item $\begin{pmatrix}8 & -8 & -8\\ 7 & -10 & -7\\ -3 & 6 & 3\end{pmatrix}$
\item $\begin{pmatrix}-2 & -8 & 4\\ 0 & 0 & 2\\ 0 & -4 & 6\end{pmatrix}$
\item $\begin{pmatrix}3 & 0 & -1\\ -2 & 1 & 1\\ 2 & 2 & 2\end{pmatrix}$
\item $\begin{pmatrix}-6 & -6 & -5\\ 0 & -2 & 0\\ 10 & 12 & 9\end{pmatrix}$
\item $\begin{pmatrix}1 & -6 & 0\\ 0 & 4 & 0\\ 2 & 5 & 3\end{pmatrix}$
\item $\begin{pmatrix}2 & -5 & 9\\ 0 & 2 & 2\\ 0 & -1 & 5\end{pmatrix}$
\item $\begin{pmatrix}0 & -2 & 2\\ 3 & 4 & -3\\ -1 & -2 & 3\end{pmatrix}$
\item $\begin{pmatrix}-2 & 2 & 3\\ 0 & -8 & -5\\ 0 & 10 & 7\end{pmatrix}$
\item $\begin{pmatrix}7 & 3 & 10\\ -4 & 0 & -4\\ -3 & -3 & -6\end{pmatrix}$
\item $\begin{pmatrix}2 & 6 & -7\\ 2 & 3 & -4\\ 2 & 4 & -5\end{pmatrix}$
\item $\begin{pmatrix}-3 & 2 & 1\\ 0 & 1 & 0\\ 0 & -6 & -2\end{pmatrix}$
\end{enumerate}

\emph{3. Further problems.}
\bigskip
\begin{enumerate}[label=\arabic*., leftmargin=*]
\item \textbf{(Characteristic polynomial.)} Find the characteristic polynomial of $A=\begin{pmatrix}1 & -2 & -2\\ -1 & 1 & 1\\ -2 & 0 & -1\end{pmatrix}$ and check how its coefficients are related to $\det A$ and $\operatorname{tr}A$.
\item \textbf{(Characteristic polynomial.)} Find the characteristic polynomial of $A=\begin{pmatrix}-2 & -1 & 0\\ 1 & 0 & 3\\ 1 & 3 & -2\end{pmatrix}$ and check how its coefficients are related to $\det A$ and $\operatorname{tr}A$.
\item \textbf{(Characteristic polynomial.)} Find the characteristic polynomial of $A=\begin{pmatrix}-1 & 2 & 3\\ 3 & 0 & -2\\ 2 & -1 & 3\end{pmatrix}$ and check how its coefficients are related to $\det A$ and $\operatorname{tr}A$.
\item \textbf{(Characteristic polynomial.)} Find the characteristic polynomial of $A=\begin{pmatrix}3 & 0 & -1\\ 3 & -1 & 0\\ 3 & 0 & 2\end{pmatrix}$ and check how its coefficients are related to $\det A$ and $\operatorname{tr}A$.
\item \textbf{(Characteristic polynomial.)} Find the characteristic polynomial of $A=\begin{pmatrix}-1 & -1 & 3\\ -1 & -1 & -1\\ 3 & 3 & 1\end{pmatrix}$ and check how its coefficients are related to $\det A$ and $\operatorname{tr}A$.
\item \textbf{(Characteristic polynomial.)} Find the characteristic polynomial of $A=\begin{pmatrix}-1 & 0 & -2\\ 2 & 1 & -2\\ -1 & 0 & -1\end{pmatrix}$ and check how its coefficients are related to $\det A$ and $\operatorname{tr}A$.
\item \textbf{(Eigenvalues of related matrices.)} The eigenvalues of $A$ are $1$, $2$, $-3$. Find the eigenvalues of $A^2$, $A^{-1}$ and $A+2E$.
\item \textbf{(Triangular matrix.)} Write down the eigenvalues of $\left(\begin{smallmatrix}2&5&-1\\0&-3&4\\0&0&1\end{smallmatrix}\right)$ without any computation.
\item \textbf{(Markov chain.)} Find the eigenvector of the transition matrix $P=\left(\begin{smallmatrix}0.9&0.2\\0.1&0.8\end{smallmatrix}\right)$ corresponding to $\lambda=1$ (the stationary distribution).
\item \textbf{(Symmetric matrix.)} Show that the eigenvectors of $\left(\begin{smallmatrix}2&1\\1&2\end{smallmatrix}\right)$ are orthogonal.
\end{enumerate}

